%%%%%%%%%%%%%%%%%%%%%%%%%%%%%%%%%%%%%%%%%%%%%%%%%%%%%%%%%%%%%%%%%%%%%%%%%%%%%%%%%%%%%
%%%%%%%%%%%%%%%%%%%%%%%%%%%%%%%%%%%%%%%%%%%%%%%%%%%%%%%%%%%%%%%%%%%%%%%%%%%%%%%%%%%%%
%%%%%%%%%%%%%%%%%%%%%%%%%%%%%%%%%%%%%%%%%%%%%%%%%%%%%%%%%%%%%%%%%%%%%%%%%%%%%%%%%%%%%
% Some comments might need to address later:

% \jian{my general comment: you mention `conceptually and technically simple' many times, be prepared to get questions like why existing works are not `conceptually' or `technically' simple, how you define `conceptually simple' and `technically simple' etc.}

% \jian{check whether these papers have ablation study, no need to mention anything here, just to be prepared if some reviewer raises questions about their ablation study}

%%%%%%%%%%%%%%%%%%%%%%%%%%%%%%%%%%%%%%%%%%%%%%%%%%%%%%%%%%%%%%%%%%%%%%%%%%%%%%%%%%%%%
%%%%%%%%%%%%%%%%%%%%%%%%%%%%%%%%%%%%%%%%%%%%%%%%%%%%%%%%%%%%%%%%%%%%%%%%%%%%%%%%%%%%%
%%%%%%%%%%%%%%%%%%%%%%%%%%%%%%%%%%%%%%%%%%%%%%%%%%%%%%%%%%%%%%%%%%%%%%%%%%%%%%%%%%%%%


% \weilin{Temporal graph learning is important problem}
In recent years, temporal graph learning has been recognized as an important machine learning problem and has become the cornerstone behind a wealth of high-impact applications~\cite{yu2018spatio,bui2021spatial,kazemi2020representation,zhou2020graph,Cong2021DyFormerAS}.
Temporal link prediction is one of the classic downstream tasks % its most important learning tasks
which focuses on predicting the future interactions among nodes.
For example, in an ads ranking system, the user-ad clicks can be modeled as a temporal bipartite graph whose nodes represent users and ads, and links are associated with timestamps indicating when users click ads. Link prediction between them can be used to predict whether a user will click an ad. % in the user-ad recommender system, the relationship between users and ads could be modeled by a temporal graph, where users and ads are nodes, and edges are created when a user clicks an ad. 
Designing graph learning models that can capture node evolutionary patterns and accurately predict future links is a crucial direction for many real-world recommender systems.


% \weilin{Existing methods are relying on RNN and SAM}

In temporal graph learning, recurrent neural network (RNN) and self-attention mechanism (SAM) have become the de facto standard for temporal graph learning~\cite{kumar2019predicting,sankar2020dysat,xu2020inductive,tgn_icml_grl2020,wang2020inductive}, and the majority of the existing works focus on designing neural architectures with one of them and additional components to learn representations from raw data.
Although powerful, these methods are conceptually
% \si{what do you mean by `conceptually' and why are they conceptually complicated?} 
and technically complicated with advanced model architectures. It is non-trivial to understand which parts of the model design truly contribute to its success, and whether these components are indispensable.
Thus, in this paper, we aim at answering the following two questions:

% \weilin{We propose ... briefly introduce our method + results}

\underline{{Q1: Are RNN and SAM always indispensable for temporal graph learning?}}
To answer this question, we propose \oure, a simple architecture based entirely on the multi-layer perceptrons (MLPs) and neighbor mean-pooling, which does not utilize any RNN or SAM in its model architecture (Section~\ref{section:method}).
% The useful information is extracted via data pre-processing and fed into \our for temporal graph link prediction. 
Despite its simplicity, \our could obtain outstanding results when comparing it against baselines that are equipped with the RNNs and SAM.
In practice, it achieves state-of-the-art performance in terms of different evaluation metrics (e.g., average precision, AUC, Recall@K, and MRR) on real-world temporal graph datasets, with the even smaller number of model parameters and hyper-parameters, and a conceptually simpler input structure and model architecture 
% \si{again, do we have any illustration to support that ours is conceptually simpler in terms of model arch?} 
(Section~\ref{section:experiments}).

% \weilin{Brief explain why model success? Need to briefly summarize our observation here ...} 

\underline{{Q2: What are the key factors that lead to the success of \oure?}}
% To achieve better understanding, we study the performance of \our in Section~\ref{section:success}. 
We identify three key factors that contribute to the success of \oure:
\circled{1} \emph{The simplicity of \oure's input data and neural architecture.}
% \si{the point of simplicity of neural architecture is already covered in the last sentence (i.e., \circled{2})}.}
Different from most deep learning methods that focus on designing conceptually complicated data preparation techniques and technically complicated neural architectures, we choose to simplifying the neural architecture and utilize a conceptually simpler data as input. Both of which could lead to a better model performance and better generalization (Section~\ref{section:data_alignment}).
\circled{2} \emph{A time-encoding function that encodes any timestamp as an easily distinguishable input vector for \oure}. Different from most of the existing methods that propose to learn the time-encoding function from the raw input data, our time-encoding function utilizes conceptually simple features and is fixed during training. Interestingly, we show that our fixed time-encoding function is more preferred than the trainable version (used by most previous studies), and could lead to a smoother optimization landscape, a faster convergence speed, and a  better generalization (Section~\ref{section:time_encoder});
\circled{3} \emph{A link-encoder that could better distinguish temporal sequences.}
Different from most existing methods that summarize sequences using SAM, our encoder module is entirely based on MLPs.
Interestingly, our encoder can distinguish temporal sequences that cannot be distinguished by SAM, 
% \si{double check if we can distinguish these sequences},
and it could generalize better due to its simpler neural architecture and lower model complexity (Section~\ref{section:expressive_power}).

% we choose to take more seriously on designing simpler neural architecture \weilin{with ... data}.
% The selection of good input data/features provides a better alignment between the input features and their labels. 
% As a result, a simple neural network model is enough to capture the underlying mapping between the inputs to their labels (Section~\ref{section:data_alignment}). 
% \weilin{existing work has complicated input data structure, we also propose to simplify the input data, which also leads to a better performance.}

To this end, we summarize our contributions as follows: 
\circled{1} We propose a conceptually and technically simple architecture \oure;
\circled{2} Even without RNN and SAM, \our not only outperforms all baselines but also enjoys a faster convergence and better generalization ability;
\circled{3} Extensive study identifies three  factors that contribute to the success of \oure. 
\circled{4} Our results could motivate future research to rethink the importance of the conceptually and technically simpler method.


